\documentclass[10pt,a4paper]{amsart}
\usepackage[margin=19mm]{geometry}
\usepackage{amsmath,amssymb,mathtools}
\usepackage[scr=rsfs,cal=euler]{mathalfa}
\usepackage{xfrac}
\usepackage[unicode,hidelinks,bookmarks=false]{hyperref}
\newcommand{\ie}{i.e.\ }
\newcommand{\cf}{cf.\ }
\newcommand{\resp}{respectively\ }
\newcommand{\Subsection}{\subsection}
\providecommand{\nobreakdash}{\nobreak\hbox{-}}
\setlength{\emergencystretch}{2em}
\multlinegap0pt
\usepackage{xcolor}
\usepackage{amssymb}
\usepackage{mathtools}
\usepackage{dsfont}
\usepackage{tikz}
\usepackage{algorithm}\usepackage{algpseudocode}
\usepackage[shortlabels]{enumitem}
\usepackage{csquotes}
\newtheorem{lemma}{Lemma}
\newcommand{\PP}{\mathbb{P}}
\title{Triangulation of Points Constrained to a Plane}
\author{Petr Hrub\'y, Elima Shehu}
\date{Source: arXiv:2604.27246v1 (CC BY 4.0)}
\begin{document}
\maketitle

\noindent\emph{Source and licence.} Triangulation of Points Constrained to a Plane. Version arXiv:2604.27246v1 (CC BY 4.0), distributed by arXiv under the Creative Commons Attribution 4.0 International licence, http://creativecommons.org/licenses/by/4.0/, as recorded in the arXiv licence metadata for this exact version and shown on the abstract page https://arxiv.org/abs/2604.27246v1. Full licence notice: \texttt{licenses/arxiv-2604.27246-CC-BY-4.0.txt}.\par\smallskip

\noindent\textit{Editorial scope.} Counted unit: the article's lemma lem:DQ-infty-plane (source0.tex lines 2136--2144). The article's complete original proof of it is delivered with it (source0.tex lines 2146--2202, one \texttt{proof} environment of 57 lines). No mathematics is written, repaired or re-derived: the statement and the proof are the article's own lines, in the article's own notation, and no retained line is substituted or re-flowed.  No further source-local context is needed: every cross-reference of every retained line resolves inside the two delivered spans.  Nothing was omitted from any retained span, so the package carries no omission record.  No retained line contains a citation, so the wrapper carries no bibliography and no bibliography entry.  No retained line contains a figure environment, an image-inclusion command or drawing code, so no image is delivered and the package declares no asset dependency.  Editorial formatting only, in addition: an AMS \texttt{amsart} preamble that restates the article's own declarations and macros used by the retained lines, namely the environments \texttt{lemma} on separate counters, together with the standard packages those lines need.  The retained environments print in this document as Lemma 1 in order of appearance, whereas the article prints the same items with the numbering of its own full document.  The complete native source of the article as distributed in the arXiv e-print for this exact version is bundled unchanged as \texttt{sources/199460/source0.tex}.
\begin{lemma}[Intersection of $D_{Q}^{\Pi}$ with $D_{\infty}^{\Pi}$]
\label{lem:DQ-infty-plane}
Let $\mathcal{M}^\Pi_{\mathcal{C}}$ be the planar-anchored point multiview variety 
for $m$ cameras in generic position with anchor plane $\Pi \subset \PP^3$.
Then
\[
\chi\!\left(D_{Q}^{\Pi}\cap D_{\infty}^{\Pi}\right) \;=\; 2m + \binom{m}{2}.
\]
\end{lemma}

\begin{proof}
Intersecting the homogenized distance divisor $H_Q$ with the hyperplane at
infinity $H_{\infty,i}=\{x_i=0\}$, one obtains three types of components
as in the general decomposition:
\[
\begin{aligned}
H_Q\cap H_{\infty,i}
&=\{y_{2i-1}+\sqrt{-1}\,y_{2i}=0,\,x_i=0\} \\
&\quad\cup\ \{y_{2i-1}-\sqrt{-1}\,y_{2i}=0,\,x_i=0\} \\
&\quad\cup\ \bigcup_{j\neq i}\{x_i=x_j=0\}.
\end{aligned}
\]

For the planar-anchored point multiview variety $\mathcal{M}^\Pi_{\mathcal C}$, set 
\[
\begin{aligned}
K_{i}^{\Pi,\pm}
&:=\mathcal{M}^\Pi_{\mathcal C}\cap\{y_{2i-1}\pm\sqrt{-1}\,y_{2i}=0,\;x_i=0\},\\
A_{i,j}^{\Pi}
&:=\mathcal{M}^\Pi_{\mathcal C}\cap\{x_i=x_j=0\},
\qquad (i\neq j).
\end{aligned}
\]
\textbf{Contribution of the $K_i^{\Pi,\pm}$ components.}
For a fixed camera $C_i$, the conditions
\[
y_{2i-1}\pm\sqrt{-1}\,y_{2i}=0,\qquad x_i=0
\]
mean that the image point in the $i$-th factor is the isotropic point
\mbox{\([0:\pm \sqrt{-1}:1]\in\PP^2\).}
Back-projecting such a point gives a \emph{complex} line through the camera center. Since the anchor plane $\Pi$ is not parallel to this direction,
this complex line meets~$\Pi$ in a unique point.
Thus each of the two isotropic directions produces a single configuration in $\mathcal{M}^\Pi_{\mathcal C}$.

Hence, for each $i$, the sets $K_i^{\Pi,+}$ and $K_i^{\Pi,-}$ each contribute
one point, giving a total of
\[
\sum_{i=1}^m \chi(K_i^{\Pi,+})+\chi(K_i^{\Pi,-})
=2m.
\]

\textbf{Contribution of the $A_{i,j}^\Pi$ components.}
The components $\{x_i=x_j=0\}$ impose that the image points of cameras $i$ and $j$ both lie at infinity. Geometrically, this means the back-projected rays of
$C_i$ and $C_j$ meet at a point of the anchor plane $\Pi$.
For generic cameras this intersection is unique, so each $A_{i,j}^\Pi$
consists of a single point, and different pairs $\{i,j\}$ give disjoint
configurations. Therefore
\[
\sum_{i<j}\chi(A_{i,j}^\Pi)=\binom{m}{2}.
\]

Adding both contributions gives
\[
\chi(D_Q^\Pi\cap D_\infty^\Pi)
=2m+\binom{m}{2}.
\]
\end{proof}

\end{document}
